% Slides of the English course "Computer Architecture" (same theme as the
% ENSB decks: light background, teal + orange). Compile from en/slides/.
\usepackage{amsmath, amssymb}
\usepackage{booktabs, colortbl, tabularx, array}
\usepackage{tikz}
\usepackage[skins]{tcolorbox}
\usepackage{fontawesome5}
\usepackage{listings}
\usepackage{fontspec}
\usetikzlibrary{positioning, calc, arrows.meta, shapes.geometric, decorations.pathreplacing}

\setsansfont{texgyreheros}[Extension=.otf, UprightFont=*-regular, BoldFont=*-bold,
    ItalicFont=*-italic, BoldItalicFont=*-bolditalic]
\newfontfamily\codefont[Scale=0.95, BoldFont=lmmonolt10-bold.otf,
    ItalicFont=lmmono10-italic.otf]{lmmono10-regular.otf}

\definecolor{brandTeal}{HTML}{5CC8E6}
\definecolor{brandOrange}{HTML}{FFA27F}
\definecolor{brandLight}{HTML}{F8F9FA}
\definecolor{brandDark}{HTML}{263238}
\definecolor{brandMuted}{HTML}{6C757D}
\definecolor{brandGreen}{HTML}{7BC47F}
\definecolor{brandPurple}{HTML}{B39DDB}
\colorlet{tealDeep}{brandTeal!60!black}
\colorlet{orangeDeep}{brandOrange!70!black}
\definecolor{codeGrey}{HTML}{F1F3F5}

\setbeamercolor{background canvas}{bg=brandLight}
\setbeamercolor{normal text}{fg=brandDark}
\setbeamercolor{structure}{fg=brandTeal}
\setbeamertemplate{navigation symbols}{}
\setbeamertemplate{frametitle}{%
    \vspace{0.35cm}
    \noindent{\bfseries\Large\color{brandDark}\insertframetitle}\par
    \vskip-2pt\noindent{\color{brandTeal}\rule{\textwidth}{1.5pt}}\vspace{0.05cm}}
\setbeamertemplate{footline}{%
    \hbox to \paperwidth{\hspace{0.6cm}{\color{brandMuted}\scriptsize\insertshorttitle}\hfill
        {\color{brandMuted}\scriptsize\insertframenumber\,/\,\inserttotalframenumber}\hspace{0.6cm}}%
    \vspace{0.25cm}}
\setbeamertemplate{itemize item}{\color{brandTeal}\small\faChevronRight}
\setbeamertemplate{itemize subitem}{\color{brandMuted}\small\faAngleRight}
\setbeamertemplate{enumerate item}{\color{tealDeep}\bfseries\insertenumlabel.}
% list lines are 1.4 cm apart (baseline to baseline)
\newlength{\lineGap}\setlength{\lineGap}{1.4cm}
\let\origlist\list
\def\list#1#2{\origlist{#1}{#2\setlength{\itemsep}{\dimexpr\lineGap-\baselineskip\relax}\setlength{\parsep}{0pt}}}

\newcolumntype{L}{>{\raggedright\arraybackslash}X}
\newcolumntype{C}{>{\centering\arraybackslash}X}
\newcommand{\ras}[1]{\cellcolor{brandTeal!20}\textbf{#1}}
\newcommand{\cell}[2]{\leavevmode\textcolor{#1}{#2}}
\renewcommand{\arraystretch}{1.3}
\setlength{\arrayrulewidth}{0.6pt}
\arrayrulecolor{brandMuted!70}

\tcbset{brandbox/.style={enhanced, arc=4pt, outer arc=4pt, boxrule=1.5pt,
    coltitle=white, coltext=brandDark, fonttitle=\bfseries,
    toptitle=0.8mm, bottomtitle=0.8mm, top=1.2mm, bottom=1.2mm, left=3mm, right=3mm,
    before skip=4pt, after skip=4pt}}
\newtcolorbox{tbox}[2][]{brandbox, colback=brandTeal!10, colframe=brandTeal, title={#2}, #1}
\newtcolorbox{obox}[2][]{brandbox, colback=brandOrange!10, colframe=brandOrange, title={#2}, #1}
\newtcolorbox{gbox}[2][]{brandbox, colback=brandGreen!8, colframe=brandGreen, title={#2}, #1}
\newtcolorbox{pbox}[2][]{brandbox, colback=brandPurple!10, colframe=brandPurple, title={#2}, #1}

\newcommand{\code}[1]{\mbox{\codefont\color{tealDeep}#1}}
\newcommand{\hex}[1]{\mbox{#1\textsubscript{H}}}
\newcommand{\bin}[1]{\mbox{#1\textsubscript{2}}}
\newcommand{\nql}[2]{\mbox{#1\,$\leftarrow$\,#2}}

% ---- code from ../../code/ (the verified programs of the course) ----
\lstdefinelanguage{asm86}{
    morekeywords={mov, add, adc, sub, sbb, inc, dec, neg, mul, imul, div, idiv,
        cmp, and, or, xor, not, test, shl, shr, sal, sar, rol, ror, jmp, je, jne,
        jz, jnz, jg, jge, jl, jle, ja, jae, jb, jbe, jcxz, loop, call, ret, iret,
        int, push, pop, lea, xchg, in, out, nop, movsb, rep, cld, std, cli, sti,
        syscall, addps, vaddps},
    morekeywords=[2]{org, db, dw, dd, dq, equ, proc, endp, offset, ptr, byte,
        word, dword, qword, section, global, dup, resb, rel},
    sensitive=false, morecomment=[l]{;}, morestring=[b]', morestring=[b]"}
\lstdefinestyle{base}{basicstyle=\codefont\fontsize{8}{8.6}\selectfont, columns=fullflexible,
    keepspaces=true, showstringspaces=false, upquote=true,
    keywordstyle=\color{tealDeep}\bfseries, keywordstyle=[2]\color{brandPurple!60!black}\bfseries,
    commentstyle=\color{red!75!black}, stringstyle=\color{orangeDeep},
    numbers=left, numberstyle=\tiny\color{brandMuted}, numbersep=7pt,
    firstnumber=auto,  % excerpts keep the real line numbers of the book
    xleftmargin=16pt, aboveskip=0pt, belowskip=0pt}
\lstdefinestyle{asm}{style=base, language=asm86}
\lstdefinestyle{c}{style=base, language=C}
\newtcolorbox{codebox}{enhanced, colback=codeGrey, colframe=codeGrey, boxrule=0pt, arc=2pt,
    left=3pt, right=3pt, top=2pt, bottom=2pt, before skip=3pt, after skip=3pt,
    overlay={\fill[brandTeal!70] (frame.south west) rectangle ([xshift=2.5pt]frame.north west);}}
\newtcolorbox{outbox}{enhanced, colback=brandDark, colframe=brandDark, boxrule=0pt, arc=2pt,
    left=5pt, right=5pt, top=2pt, bottom=2pt, before skip=3pt, after skip=3pt, colupper=white}
% \asm[firstline=5, lastline=20]{ch3/procnum.asm}   \cfile{...}   \screen{ch3/hello.out}
% assembly on slides: copies in code/ with inline comments removed (same lines,
% same numbers as the book; the explanation is on the slide)
\newcommand{\asm}[2][]{\begin{codebox}\lstinputlisting[style=asm,#1]{code/#2}\end{codebox}}
% smaller code for programs with long comments
\newcommand{\codesize}{\codefont\fontsize{7}{7.5}\selectfont}
\newcommand{\asms}[2][]{\begin{codebox}\lstinputlisting[style=asm,basicstyle=\codesize,#1]{code/#2}\end{codebox}}
\newcommand{\cfiles}[2][]{\begin{codebox}\lstinputlisting[style=c,basicstyle=\codesize,#1]{../../code/#2}\end{codebox}}
\newcommand{\cfile}[2][]{\begin{codebox}\lstinputlisting[style=c,#1]{../../code/#2}\end{codebox}}
\newcommand{\screen}[2][]{\begin{outbox}\lstinputlisting[basicstyle=\codefont\fontsize{8}{8.6}\selectfont\color{white},
    columns=fullflexible, keepspaces=true, aboveskip=0pt, belowskip=0pt,#1]{../../code/#2}\end{outbox}}

\tikzset{
    sahm/.style={-{Stealth[length=2.4mm]}, thick, color=brandMuted},
    sahmT/.style={-{Stealth[length=2.4mm]}, very thick, color=tealDeep},
    sahmO/.style={-{Stealth[length=2.4mm]}, very thick, color=orangeDeep},
    unit/.style={draw=#1, fill=#1!15, very thick, rounded corners=4pt,
        align=center, font=\bfseries\small, inner sep=4pt, text=brandDark},
    unit/.default=brandTeal,
    reg/.style={draw=brandDark, fill=white, thick, minimum height=0.6cm,
        font=\codefont\small, inner sep=2pt},
    % styles with a colour argument must be defined here, not inside a frame
    layer/.style={draw=#1, fill=#1!15, very thick, minimum width=9cm, minimum height=0.7cm, font=\small},
    evo/.style={unit=#1, minimum width=2.6cm, minimum height=1.2cm, font=\small},
}

% Title slide: \cover{Chapter 1}{Title}{subtitle}
\newcommand{\cover}[3]{%
\begin{frame}[plain]
    \begin{tikzpicture}[remember picture, overlay]
        \draw[brandTeal, line width=1.5pt, opacity=0.3] (current page.south east) ++(-1,1) circle (4cm);
        \draw[brandTeal, line width=0.8pt, opacity=0.25] (current page.north west) ++(2,2) circle (6cm);
        \draw[brandOrange, line width=0.5pt, opacity=0.25] (current page.north west) ++(2,2) circle (7cm);
        \begin{scope}[shift={([xshift=-1.6cm, yshift=-1.4cm]current page.north east)}, scale=0.35]
            \fill[brandDark] (-1.5,-1.5) rectangle (1.5,1.5);
            \fill[brandTeal] (-1,-1) rectangle (1,1);
            \foreach \i in {-1.2,-0.6,0,0.6,1.2} {
                \draw[brandMuted, line width=1pt] (\i,1.5) -- (\i,2); \draw[brandMuted, line width=1pt] (\i,-1.5) -- (\i,-2);
                \draw[brandMuted, line width=1pt] (1.5,\i) -- (2,\i); \draw[brandMuted, line width=1pt] (-1.5,\i) -- (-2,\i);}
        \end{scope}
    \end{tikzpicture}
    \vspace{1.4cm}
    \begin{minipage}{0.85\textwidth}
        {\color{tealDeep}\bfseries Computer Architecture \textbullet\ #1}\\[0.3cm]
        {\color{brandDark}\LARGE\bfseries #2\par}\vspace{0.3cm}
        {\color{brandMuted}\large #3}
    \end{minipage}
    \vspace{1.2cm}
    \begin{minipage}{0.85\textwidth}
        \color{brandMuted}\small
        \begin{tabular}{@{}ll}
            \textbf{Teacher:} & Dr. Benabderrezak Youcef \\
            \textbf{Course:} & UEF31, Architecture des ordinateurs \\
        \end{tabular}
    \end{minipage}
\end{frame}}
